
\documentclass[12pt]{amsart}

\usepackage[margin=1in]{geometry}
\usepackage{amsfonts}
\usepackage{amsthm}
\usepackage{thmtools}
\usepackage{amsmath}
\usepackage{amssymb}
\usepackage{graphicx}
\usepackage{mathrsfs}
\usepackage{array}
\usepackage{amssymb}
% \usepackage{mathabx}
\usepackage{tikz-cd}
\usepackage{mathtools}
\usepackage{hyperref}
\usepackage[nameinlink, capitalize]{cleveref}
\hypersetup{
    colorlinks=true,
    linkcolor=black,
    filecolor=green,      
    urlcolor=cyan,
    citecolor=magenta,
    pdftitle={Notes of my current research},
    pdfpagemode=FullScreen,
}
\urlstyle{same}
% \usepackage{caption} 
% \captionsetup[table]{skip=10pt}
% \usepackage{subfiles}



\newcommand{\R}{\mathbb{R}}
\newcommand{\Q}{\mathbb{Q}}
\newcommand{\V}{\mathbb{V}}
\newcommand{\sph}{\mathbb{S}}
\newcommand{\C}{\mathbb{C}}
\newcommand{\Z}{\mathbb{Z}}
\newcommand{\End}{\mathrm{End}}
\newcommand{\Aut}{\mathrm{Aut}}
\newcommand{\Hom}{\mathrm{Hom}}
\newcommand{\id}{\mathrm{id}}

\newcommand{\bbA}{\mathbb{A}}
\newcommand{\bbH}{\mathbb{H}}
\newcommand{\bbP}{\mathbb{P}}

\newcommand{\cM}{\mathcal{M}}
\newcommand{\cL}{\mathcal{L}}
\newcommand{\cE}{\mathcal{E}}
\newcommand{\cT}{\mathcal{T}}
\newcommand{\cV}{\mathcal{V}}
\newcommand{\cW}{\mathcal{W}}
\newcommand{\cH}{\mathcal{H}}
\newcommand{\cP}{\mathcal{P}}
\newcommand{\cA}{\mathcal{A}}
\newcommand{\cG}{\mathcal{G}}
\newcommand{\cU}{\mathcal{U}}

\newcommand{\IH}{\mathrm{IH}}

\newcommand{\Sp}{\mathrm{Sp}}
\newcommand{\SL}{\mathrm{SL}}
\newcommand{\GL}{\mathrm{GL}}
\newcommand{\U}{\mathrm{U}}
\newcommand{\gr}{\mathrm{gr}}
\newcommand{\ol}{\overline}
\newcommand{\ul}{\underline}
\newcommand{\wt}{\widetilde}
\newcommand{\wh}{\widehat}
\newcommand{\codim}{\mathrm{codim}}

\newcommand{\rightQ}[2]{\left.\raisebox{.2em}{$#1$}\middle/\raisebox{-.2em}{$#2$}\right.}
\newcommand{\leftQ}[2]{\left.\raisebox{-.2em}{$#2$}\middle\backslash\raisebox{.2em}{$#1$}\right.}

\DeclareMathOperator{\im}{\mathrm{im}}


\NewDocumentCommand{\ltranspose}{m}{
    \ensuremath{\prescript{\mathsf{T}}{}{{#1}}}
}

\NewDocumentCommand{\transpose}{m}{
    \ensuremath{{#1}^{\mathsf{T}}}
}


\newtheorem{theorem}{Theorem}[section]
\newtheorem{proposition}[theorem]{Proposition}
\newtheorem{lemma}[theorem]{Lemma}
\newtheorem{corollary}[theorem]{Corollary}
\newtheorem{question}[theorem]{Question}
\newtheorem{conjecture}[theorem]{Conjecture}

\theoremstyle{definition}
\newtheorem{definition}[theorem]{Definition}
\newtheorem{example}[theorem]{Example}
\newtheorem{remark}[theorem]{Remark}

\numberwithin{equation}{section}




\title{Research statement}
\author{Edward Varvak}
\address{\parbox{\linewidth}{
Department of Mathematics \\
University of Chicago \\
5734 S University Ave \\
Chicago, IL 60637 \\}}

\email{varvak@uchicago.edu}


\begin{document}

\maketitle

\section{Summary of my results}

My research concerns the structure of Lagrangian fibration on symplectic varieties. 

\begin{definition} A normal variety \(X\) is {\bf symplectic} if it admits a non-degenerate 
        symplectic form \(\sigma \in H^0(X^\mathrm{sm}, \Omega_{X^\mathrm{sm}}^2)\) which 
        extends holomorphically onto a resolution of singularities \(\pi: Y \to X\).
        In this setting, a {\bf Lagrangian fibration} is a proper morphism \(f: X \to B\) 
        with connected half-dimensional fibers, such that \(\sigma|_{X_b} = 0\) for any 
        fiber \(X_b\).
\end{definition}

We will need to assume that \(X\) is compact for most of our results. Additionally, 
we will want \(X\) to be ``irreducible'' in a sense, which will be made more precise 
in (*** link below). To that end, the most general class of symplectic varieties we 
will work with is the primitive symplectic varieties. 

\begin{definition}
    A compact symplectic variety \(X\) is {\bf primitive symplectic} if it satisfies 
    \(H^1(X, \mathcal{O}_X) = 0\) and \(\sigma\) is unique up to scaling, i.e. 
    \(H^0(X, \Omega_X^2) = \C \sigma\).
\end{definition}

As we will see (*** link below), when \(X\) is a primitive symplectic variety admitting 
a Lagrangian fibration \(f: X \to B\), the smooth fibers \(X_b\) are Abelian varieties. 
My main results further classify the structure of the fibers.

\begin{theorem} \label{mt-classification}
    Let \(X\) be a primitive symplectic variety of dimension \(2n\) admitting a 
    Lagrangian fibration \(f: X \to B\). If \(f\) is not isotrivial, the Mumford--Tate 
    group of the very general fiber \(X_b\) is isogenous to \((\Sp_{2h})^{n/h}\).
\end{theorem}

Recall that an Abelian variety is Hodge-generic if it admits no extra Hodge tensors. 
In this case, it is simple and carries no extra symmetries.

\begin{corollary}
    The very general fiber \(X_b\) is isogenous to a product of Abelian varieties 
    \(A_1 \times \cdots \times A_{n/h}\), where each \(A_i\) is Hodge-generic and 
    \(A_i\) is not isogenous to \(A_j\) for \(i \neq j\).
\end{corollary}

The most well-studied class of primitive symplectic varieties consists of hyper-K\"ahler 
manifolds, which should be thought of as the natural generalization of K3 surfaces.

\begin{definition}
    A compact complex manifold \(X\) is {\bf hyper-K\"ahler} if \(X\) is simply connected 
    and its algebra of global holomorphic forms is generated by a non-degenerate 
    symplectic form: 
    \(H^0(X, \Omega_X^{2p}) = \C \sigma^p,\quad H^0(X, \Omega_X^{2p+1}) = 0.\)
\end{definition}

% In the case where \(X\) is a sufficiently general among hyper-K\"ahler manifolds 
% admitting a Lagrangian fibration, we can say much more about the Mumford--Tate group 
% of its very general fiber. 

\begin{theorem}
    Let \(\mathcal{X}\) be a deformation class of hyper-K\"ahler manifolds such that 
    \begin{enumerate}
        \item all members \(X \in \mathcal{X}\) have second Betti number 
        \(b_2(X) \geq 7\),
        \item there exists some \(X \in \mathcal{X}\) admitting a Lagrangian fibration.
    \end{enumerate}
    Choose \(X\) to be general among the hyper-K\"ahler manifolds admitting a Lagrangian 
    fibration, \(f: X \to B\). If \(X_b\) is a very general fiber of \(f\), 
    then it has Mumford--Tate group \(\mathbf{MT}(X_b) = \mathrm{GSp}_{2n}\), i.e. 
    \(X_b\) is Hodge-generic.
\end{theorem}

These results largely came about through understanding the monodromy of the family \(f\). 
To be more precise, let \(f: X^\circ \to B^\circ\) be the restriction of \(f\) to 
its regular points, and let \(V = R^1 f_* \Z_{X^\circ}\), viewed as a local system. 

\begin{theorem}
    Suppose that \(X\) is primitive symplectic, \(f: X \to B\) a Lagrangian fibration, 
    \(V_\C = V \otimes_\Z \C\) the underlying \(\C\)-representation of 
    \(\pi_1(B^\circ)\). 
    \begin{enumerate}
        \item If \(f\) is not isotrivial, then \(V_\C\) is irreducible.
        \item If \(f\) is isotrivial, then its smooth fibers are isogenous to 
        \(E^n\), for \(E\) an elliptic curve. Furthermore, if \(K\) is the endomorphism 
        field of \(E\) (possibly \(\Q\)), then \(V_K = U_1 \oplus U_2\) splits into 
        two geometrically irreducible components.
    \end{enumerate}
\end{theorem}

\begin{remark}
There are a few clarifications to be made regarding the results and contributions. 
For the non-isotrivial case (1), the later result of \Cref{mt-classification} 
determines that a finite-index subgroup of the image of \(\pi_1(B^\circ)\) lies in 
\((\mathrm{Sp}_{2h, \Q})^{n/h}\), and acts on \(V_\C\) via a direct sum of standard representations. 
Such a representation is would not be irreducible, if not for the unaccounted finite 
part of the monodromy transitively permuting the irreducible factors.
For the isotrivial case (ii), the result \(X_b \sim E^n\) was first proven in 
(*** KLM), though we produced an alternative proof of this fact. 
\end{remark}




\section{Background and motivation}

Much of modern research is directed at understanding the deformation classes of 
hyper-K\"ahler manifolds. There are currently four known deformation types: the 
infinite families \(K3^{[n]}\) and \(\mathrm{Kum}_n\), and the two sporadic types 
\(OG6\) and \(OG10\) (loc. cit. O'Grady). The SYZ Conjecture (cite below) predicts 
that every deformation class \(\mathcal{X}\) 






\section{Future directions}



\end{document}